\documentclass[11pt]{article}
\usepackage[a4paper,margin=1in]{geometry}
\usepackage{amsmath,amssymb,mathtools,bm}
\usepackage{enumitem,booktabs,array}
\usepackage{tcolorbox}
\usepackage{hyperref}
\usepackage[protrusion=true,expansion=false]{microtype}
\hypersetup{colorlinks=true,linkcolor=blue,urlcolor=blue,citecolor=blue}
\setlist[itemize]{topsep=2pt,itemsep=2pt}
\setlist[enumerate]{topsep=2pt,itemsep=2pt}
\newtcolorbox{keybox}{colback=blue!3!white,colframe=blue!40!black,boxrule=0.4pt,arc=1mm}
\newtcolorbox{alertbox}{colback=red!3!white,colframe=red!40!black,boxrule=0.4pt,arc=1mm}
\newtcolorbox{answerbox}{colback=green!3!white,colframe=green!40!black,boxrule=0.4pt,arc=1mm}
\newtcolorbox{drillbox}{colback=yellow!5!white,colframe=orange!60!black,boxrule=0.4pt,arc=1mm}
\newcommand{\EV}{\mathbb{E}}
\newcommand{\Var}{\mathrm{Var}}
\newcommand{\Cov}{\mathrm{Cov}}
\newcommand{\blank}[1]{\underline{\hspace{#1}}}

\title{E03-03: Eckbo, Nygaard \& Thorburn (2022, MS) \\
\large ``Valuation effects of Norway's board gender-quota law revisited'' --- Active Recall Workbook}
\author{Empirical Corporate Finance II --- Exam Prep}
\date{\today}
\begin{document}\maketitle

\begin{keybox}
\textbf{Paper in one sentence.} Using an expanded sample and corrected identification (proper control group, longer pre- and post-windows, industry matching), ENT find that Norway's 2003--2008 gender quota caused essentially \emph{no} negative valuation effect on treated ASA firms --- overturning Ahern-Dittmar's (2012) headline finding and suggesting that board-diversity mandates need not be value-destroying.

\textbf{Placement.} Direct reassessment of Ahern-Dittmar (E03-02). Jointly they define the empirical debate on diversity mandates: AD finds large negative effects with shorter windows and noisier control, ENT finds near-zero effects with improved design.
\end{keybox}

\section*{Round 1: Complete Walk-Through}

\subsection*{1.1 Critique of AD}
\begin{itemize}
\item AD's control group (AS firms) differs on observables and is small.
\item Short post-window mixes quota effects with GFC (2008 event).
\item AD compares Tobin's Q at firms that differ in ways unrelated to the quota.
\end{itemize}

\subsection*{1.2 ENT design improvements}
\begin{itemize}
\item Larger firm sample with richer covariates.
\item Matched DiD: ASA firms matched to AS on industry, size, profitability pre-2003.
\item Longer post-window (through 2013) separating quota implementation from GFC effects.
\item Includes both book-value leverage and alternative performance measures.
\end{itemize}

\subsection*{1.3 Headline results}
\begin{itemize}
\item Matched-DiD estimate of Tobin's Q: statistically indistinguishable from zero.
\item Operating performance (ROA, ROS) effects also null.
\item AD's $-3.5$\% announcement return reflects pre-quota differences rather than a causal effect.
\item Board composition and firm size evolved, but performance did not decline.
\end{itemize}

\subsection*{1.4 Mechanism reconsidered}
\begin{itemize}
\item New female directors, while younger on average, were not less effective (measured via attendance, committee assignments).
\item Talent-pool depth rose over time, possibly endogenously to the quota.
\item The ``forced-search'' friction AD postulated is small or transient.
\end{itemize}

\subsection*{1.5 Implications}
\begin{itemize}
\item Diversity mandates can be implemented without value destruction if design is careful.
\item Policy-relevant: countries considering similar quotas (e.g.\ California, EU) need not expect large negative effects.
\item Methodological reminder on the importance of matched controls in DiD.
\end{itemize}

\subsection*{1.6 Caveats}
\begin{alertbox}
\begin{itemize}
\item Null results can still reflect power limitations; ENT argues sample is large enough.
\item Long-run cultural effects (mentorship, pipeline) are not captured by Q.
\item Estimates specific to Norway's institutional setting.
\end{itemize}
\end{alertbox}

\section*{Round 2: Level 1 Fill-in}
\begin{enumerate}
\item Paper that ENT revisits: Ahern \& \blank{1.5cm} (2012).
\item Critique of AD: inadequate \blank{2cm} group and short post-window.
\item ENT design: \blank{1.5cm} DiD with industry-size-profitability matching.
\item Main finding on Q: \blank{1.5cm} from zero.
\item Policy lesson: quotas \blank{1cm} inherently value-destroying.
\end{enumerate}

\section*{Round 3: Level 2 Prompts}
\begin{enumerate}
\item Summarize ENT's critique of Ahern-Dittmar's identification.
\item Describe the matched-DiD design and its advantages.
\item Report the main findings on Q and operating performance.
\item Discuss policy implications for gender-diversity mandates.
\item How should the AD-ENT debate be synthesized?
\end{enumerate}

\section*{Round 4: Minimal Scaffolding}
From a blank page: (i) AD critique, (ii) ENT design, (iii) null result, (iv) mechanism reconsidered, (v) policy.

\section*{Round 5: Blank Page}
Essay: what do we really know about the valuation effects of board-diversity quotas?

\vspace{6pt}
\begin{drillbox}
\textbf{Speed Drill.} (1) Paper revisited. (2) Design improvement. (3) Main Q finding. (4) Mechanism update. (5) Policy lesson.
\end{drillbox}

\section*{Quick Reference \& Answer Key}
\begin{answerbox}
\begin{itemize}
\item \textbf{Critique.} AD control group flawed; window too short.
\item \textbf{Design.} Matched DiD, longer window.
\item \textbf{Finding.} Q effect indistinguishable from zero.
\item \textbf{Lesson.} Quotas need not destroy value.
\end{itemize}
\end{answerbox}

\begin{keybox}
\textbf{30-second exam pitch.} Eckbo, Nygaard \& Thorburn (2022) revisit Ahern-Dittmar (2012) on Norway's 40\% female board-representation quota using improved identification: a larger sample, matched difference-in-differences on industry, size, and pre-2003 profitability, and a longer post-window (through 2013) that separates the quota from GFC effects. They find that Tobin's Q and operating-performance effects of the quota are statistically indistinguishable from zero, overturning AD's $-3.5$\% announcement return and $\sim$12--15\% Q decline. Newly-appointed female directors, although younger, were not less effective in attendance or committee work. The paper concludes that diversity mandates need not be value-destroying and cautions against generalizing from earlier short-window studies. Jointly with AD, it is the canonical debate on gender-quota valuation effects and a methodological lesson on the importance of matched controls in DiD.
\end{keybox}

\end{document}
